% Zone „Das kennst du schon“ – aus bank/extremalprobleme/zone.jsonl (zusammenbau v0.1)
\einheitenkopf[zone][Kennst du schon]{Das kennst du schon}
\setcounter{aufgabe}{0}

%% TODO Titel: Ich-kann-Satz fehlt in der Bank (Platzhalter: Kettenname) – Nr. 1
%% TODO Anweisung: ein Satz über der ersten Teilaufgabe fehlt in der Bank – Nr. 1
\begin{aufgabe}{Flächen- und Umfangsformeln von Rechteck, Dreieck und Trapez, Satz des Pythagoras für Streckenlängen}
\begin{teile}
\teil Berechne den Flächeninhalt eines Dreiecks mit der Grundseite 6 cm und der Höhe 4 cm. \leerfeld[cm²]
\teil Ein Trapez hat die parallelen Seiten $5{,}4$ cm und 3 cm und die Höhe $2{,}5$ cm. Berechne seinen Flächeninhalt. \leerfeld[cm²]
\end{teile}
\end{aufgabe}

%% TODO Titel: Ich-kann-Satz fehlt in der Bank (Platzhalter: Kettenname) – Nr. 2
%% TODO Anweisung: ein Satz über der ersten Teilaufgabe fehlt in der Bank – Nr. 2
\begin{aufgabe}{Punkte (x | f(x)) im Koordinatensystem lesen und Figuren einzeichnen, Funktionswerte berechnen}
\begin{teile}
\teil Berechne $f(3)$ für $f(x) = x^2 - 2x$ und gib den Punkt des Graphen an der Stelle 3 an.
\teil Berechne $f(-1{,}5)$ für $f(x) = 2x^2 + x$.
\end{teile}
\end{aufgabe}

%% TODO Titel: Ich-kann-Satz fehlt in der Bank (Platzhalter: Kettenname) – Nr. 3
%% TODO Anweisung: ein Satz über der ersten Teilaufgabe fehlt in der Bank – Nr. 3
\begin{aufgabe}{Terme ausmultiplizieren und zusammenfassen}
\begin{teile}
\teil Multipliziere aus: $x \cdot (10 - x)$
\teil Multipliziere aus: $\tfrac{1}{2}x \cdot (6 - 2x^2)$
\end{teile}
\end{aufgabe}

%% TODO Titel: Ich-kann-Satz fehlt in der Bank (Platzhalter: Kettenname) – Nr. 4
%% TODO Anweisung: ein Satz über der ersten Teilaufgabe fehlt in der Bank – Nr. 4
\begin{aufgabe}{Ableitungen bilden: Potenz-, Faktor- und Summenregel; GK und LK zusätzlich Produkt- und Kettenregel für e-Funktions-Produkte}
\begin{teile}
\teil Leite ab: $f(x) = 3x^2 - 5x$ f'(x) = \leerfeld
\teil Leite ab: $h(x) = -0{,}5x^4 + \tfrac{2}{3}x^3$ h'(x) = \leerfeld
\end{teile}
\end{aufgabe}

%% TODO Titel: Ich-kann-Satz fehlt in der Bank (Platzhalter: Kettenname) – Nr. 5
%% TODO Anweisung: ein Satz über der ersten Teilaufgabe fehlt in der Bank – Nr. 5
\begin{aufgabe}{Gleichungen lösen: linear, quadratisch mit Lösungsformel, biquadratisch mit Substitution, Ausklammern und Nullprodukt}
\begin{gleichungsraster}[2]
\gl{\text{Löse $x^2 - 9 = 0$.}} &
\gl{\text{Löse $2x^2 - 3x - 2 = 0$.}} \\
\end{gleichungsraster}
\end{aufgabe}

%% TODO Titel: Ich-kann-Satz fehlt in der Bank (Platzhalter: Kettenname) – Nr. 6
%% TODO Anweisung: ein Satz über der ersten Teilaufgabe fehlt in der Bank – Nr. 6
\begin{aufgabe}{Extrempunkte über notwendige und hinreichende Bedingung, Randvergleich auf einem Intervall}
\begin{teile}
\teil Bestimme den Hochpunkt von $f(x) = -x^2 + 6x$. H( \leerfeld | \leerfeld )
\teil Bestimme die Extrempunkte von $f(x) = x^3 - 3x$.
\end{teile}
\end{aufgabe}

%% TODO Titel: Ich-kann-Satz fehlt in der Bank (Platzhalter: Kettenname) – Nr. 7
%% TODO Anweisung: ein Satz über der ersten Teilaufgabe fehlt in der Bank – Nr. 7
\begin{aufgabe}{Flächen- und Umfangsformeln von Rechteck, Dreieck und Trapez, Satz des Pythagoras für Streckenlängen – Fehler finden}
\begin{teile}
\teil Finde den Fehler. Ein Dreieck hat die Grundseite 8 cm und die Höhe 5 cm. Rechnung: $A = 8 \cdot 5 = 40$ cm².
\end{teile}
\end{aufgabe}

%% TODO Titel: Ich-kann-Satz fehlt in der Bank (Platzhalter: Kettenname) – Nr. 8
%% TODO Anweisung: ein Satz über der ersten Teilaufgabe fehlt in der Bank – Nr. 8
\begin{aufgabe}{Flächen- und Umfangsformeln von Rechteck, Dreieck und Trapez, Satz des Pythagoras für Streckenlängen – selbst rechnen}
\begin{teile}
\teil Ein rechtwinkliges Dreieck hat die Katheten 9 cm und 7 cm. Berechne seinen Flächeninhalt. \leerfeld[cm²]
\end{teile}
\end{aufgabe}
